\begingroup
\singlespacing
\chapter{Hyperparameters and Implementation Details}
\label{app:hyperparams}

\section{Software}

All experiments use Python 3.12, PyTorch 2.x (CUDA 12.1), stable-baselines3 2.3.2, sb3-contrib 2.3.2, Gymnasium 0.29+, and MuJoCo 3.7.
Spectral filter computation uses \texttt{torch.linalg.eigh} in float64; all other computation is float32.

\section{Compute}

Experiments ran on Princeton Della (A100 GPUs) and Adroit (A100 MIG 20GB) HPC clusters.
Teacher training: 2--4 hours per 2M steps.
Student distillation: 5--30 minutes per model.
Total compute: approximately 200 A100 GPU-hours.

\section{Teacher Training}

Teachers were trained using RecurrentPPO from the sb3-contrib library with the hyperparameters in Table~A.1.
Each teacher was trained for 2M timesteps with checkpoints saved at 200K, 500K, 1M, and 2M steps.
Teacher evaluation used 20 deterministic rollouts per checkpoint.
All MuJoCo teachers used a single seed (42) for the distillation sweep; a second seed (123) was trained but not used in the main results, which is acknowledged as a limitation.

\begin{table}[ht]
\centering
\caption{RecurrentPPO teacher hyperparameters.}
\begin{tabular}{ll}
\toprule
Parameter & Value \\
\midrule
Algorithm & RecurrentPPO (sb3-contrib) \\
Policy & MlpLstmPolicy \\
LSTM hidden size & 128 \\
Total timesteps & 2{,}000{,}000 \\
Learning rate & $3 \times 10^{-4}$ \\
Batch size & 64 \\
n\_steps & 2048 \\
n\_epochs & 10 \\
$\gamma$ & 0.99 \\
GAE $\lambda$ & 0.95 \\
Clip range & 0.2 \\
Seeds & 42, 123 \\
Checkpoints & 200K, 500K, 1M, 2M steps \\
\bottomrule
\end{tabular}
\end{table}

\section{Student Distillation}

Student training used behavioral cloning with MSE loss on actions, with the hyperparameters in Table~A.2.
Each student was trained for up to 100 epochs with early stopping (patience 10 on validation loss).
Demonstrations were chunked into overlapping windows of 256 timesteps for batched training.
Evaluation used 20 rollouts in the POMDP environment per seed.

\begin{table}[ht]
\centering
\caption{Student distillation hyperparameters.}
\begin{tabular}{ll}
\toprule
Parameter & Value \\
\midrule
$d_\text{model}$ & 64 (312K-matched: 160) \\
Number of layers & 2 \\
Spectral filters (STU) & 16 \\
Optimizer & AdamW \\
Weight decay & 0.01 \\
Learning rate & $3 \times 10^{-4}$ \\
Schedule & Cosine annealing \\
Epochs & 100 (early stopping, patience 10) \\
Batch size & 64 \\
Sequence window & 256 steps \\
Gradient clipping & 1.0 \\
Demonstration episodes & 500 \\
Train/val split & 80/20 \\
Evaluation episodes & 20 \\
Seeds & 42, 123, 456 \\
\bottomrule
\end{tabular}
\end{table}

\section{Student Architecture Details}

Table~A.3 provides detailed parameter counts, model sizes, and inference latency for all architectures at $d_\text{model}=64$ on HalfCheetah-POMDP (obs=8, act=6).
Latency was measured on an Apple M2 Max CPU for single-timestep inference.
For deployment on actual edge hardware (e.g., Jetson Nano, Raspberry Pi), latencies would differ; the M2 Max numbers serve as a standardized reference for relative comparison.

\begin{table}[ht]
\centering
\caption{Detailed student model comparison at $d_\text{model}=64$, 2 layers (HalfCheetah-POMDP, obs=8, act=6). Latency measured on Apple M2 Max CPU, single-timestep inference.}
\begin{tabular}{lrrrr}
\toprule
Architecture & Parameters & Size (KB) & Latency (ms) & Throughput (steps/s) \\
\midrule
MLP           & 5{,}126 & 20 & 0.10 & 10{,}000 \\
FrameStack    & 8{,}710 & 34 & 0.20 & 5{,}000 \\
GRU           & 50{,}886 & 199 & 1.26 & 796 \\
LSTM          & 67{,}526 & 264 & 1.33 & 751 \\
GRU$_{160}$   & 311{,}526 & 1{,}217 & 2.80 & 357 \\
Transformer   & 232{,}006 & 906 & 3.36 & 297 \\
STU (full)    & 338{,}118 & 1{,}321 & 4.38 & 228 \\
\bottomrule
\end{tabular}
\end{table}

\section{Spectral Filter Computation}

The STU's fixed convolutional filters are derived from the top eigenvectors of a Hankel matrix, following \citet{agarwal2024spectral}.
The Hankel matrix is defined as $Z_{ij} = \frac{2}{(i+j)^3 - (i+j)}$ for $i,j \in \{1, \ldots, L\}$.
We extract the top-$K$ eigenvectors via \texttt{torch.linalg.eigh} in float64 (ascending eigenvalues, reversed) and scale each by $\lambda_k^{1/4}$.
When the sequence length $L < K$, we clamp $K = \min(K, L)$.

\section{Data Efficiency Experiments}

To test how quickly each architecture learns from limited demonstrations, we sweep over demonstration budgets of $\{5, 10, 25, 50, 100, 250, 1000\}$ episodes on both CopyTask ($N{=}10$) and T-Maze ($L{=}100$), with 3 seeds per condition.
Results are reported in \cref{fig:synthetic-summary}c and the accompanying text in \cref{sec:synthetic}.

\section{Code and Reproducibility}

All code, experiment scripts, and trained checkpoints are available at \url{https://github.com/wasumayan/wasumind}.
\endgroup
